\documentclass[10pt,a4paper]{amsart}
\usepackage[margin=19mm]{geometry}
\usepackage{amsmath,amssymb,mathtools}
\usepackage[scr=rsfs,cal=euler]{mathalfa}
\usepackage{xfrac}
\usepackage[unicode,hidelinks,bookmarks=false]{hyperref}
\newcommand{\ie}{i.e.\ }
\newcommand{\cf}{cf.\ }
\newcommand{\resp}{respectively\ }
\newcommand{\Subsection}{\subsection}
\providecommand{\nobreakdash}{\nobreak\hbox{-}}
\setlength{\emergencystretch}{2em}
\multlinegap0pt
\usepackage{bm}
\usepackage{bbm}
\usepackage{dsfont}
\usepackage{xspace}
\usepackage{cancel}
\usepackage{mathtools}
\usepackage{amsthm}
\makeatletter
\@ifundefined{thm}{\newtheorem{thm}{Thm}}{}
\@ifundefined{lemma}{\newtheorem{lemma}[thm]{Lemma}}{}
\makeatother
\makeatletter
\expandafter\ifx\csname abstract\endcsname\relax
\renewenvironment{abstract}
  {{\bfseries\noindent{\abstractname}\par\nobreak}\footnotesize}
  {\bigskip}
\fi
\makeatother
\title[Spectrum for a non-unitary one-dimensional two-state quantum]{Spectrum for a non-unitary one-dimensional two-state quantum walk with one defect}
\author{Takako Endo, Yohei Matsumoto, Hiromichi Ohno, Akito Suzuki}
\date{Source: arXiv:2511.04098v1 (CC BY 4.0)}
\begin{document}
\maketitle

\noindent\emph{Source and licence.} Spectrum for a non-unitary one-dimensional two-state quantum walk with one defect. Version arXiv:2511.04098v1 (CC BY 4.0), distributed under the Creative Commons Attribution 4.0 International licence, as recorded in the arXiv licence metadata for this exact version and shown on the abstract page https://arxiv.org/abs/2511.04098v1. Full rights notice: licenses/arxiv-2511.04098-CC-BY-4.0.txt.\par\smallskip

\noindent\textit{Editorial scope.} Counted unit: the Lemma whose source label is \texttt{lemma47} in the article (source0.tex lines 741--746, source label \texttt{lemma47}), delivered together with the article's own complete original proof of it (source0.tex lines 747--935, one \texttt{proof} environment of 189 lines). No mathematics is written, repaired or re-derived: statement and proof are the article's own text line for line. Required context retained uncounted: lemma46 (lines 725--736). Every definition, display and label of those lines is retained unchanged. Omitted from this package and not redistributed: 1--724, 737--740, 936--1204. Nothing inside a retained excerpt is omitted or altered: every retained body is delivered exactly as the article's lines are written, so no omission receipt of any kind accompanies this package. Citations: the retained excerpts carry no citation command at all, so no bibliography entry of the article is transcribed and no reference is redistributed. Reference adaptation: none was needed and none was made; every reference inside the delivered unit resolves inside it, so no unresolved reference or citation is produced. Printed slips kept literal: no printed word, symbol, hypothesis or number of the counted statement or of its proof was altered. Layout and class: the article's own class is \texttt{article} with packages including inputenc, babel, amsmath, amssymb, amsfonts, amsthm, latexsym, bm, textcomp, graphicx, xcolor, hyperref, grffile, bmpsize; the wrapper is the lane's pinned \texttt{amsart} editorial wrapper at 10pt on A4, which restates the theorem-like environments the retained text needs and adds the article's own macro definitions that the retained text needs (none), with extra wrapper packages (bm, bbm, dsfont, xspace, cancel, mathtools). The delivered numbering is the unit's own. Assets: no retained span contains a figure environment, a graphics-inclusion command, a PDF-page inclusion, a table or a TikZ picture; no image file is copied into this package, the package declares no asset dependency and no image file is redistributed with it. Licence: version arXiv:2511.04098v1 of this article is distributed under the Creative Commons Attribution 4.0 International licence, as recorded in the arXiv licence metadata for this exact version and shown on the abstract page https://arxiv.org/abs/2511.04098v1. A full rights notice is delivered at licenses/arxiv-2511.04098-CC-BY-4.0.txt. Characters: every retained excerpt is pure ASCII, so the delivered engine, XeLaTeX with its default Latin Modern text font, reports no missing character.
\begin{lemma}\label{lemma46}
For an arbitrary complex number $a+bi \ (a,b\in \mathbb{R})$, its square root $\sqrt{a+bi}$ is express as 

\begin{eqnarray}
\sqrt{a+bi} = \left\{
\begin{array}{ll}
\sqrt{\frac{a+\sqrt{a^2+b^2}}{2}}+i\sqrt{\frac{-a+\sqrt{a^2+b^2}}{2}} &  b\geq0,\\
\sqrt{\frac{a+\sqrt{a^2+b^2}}{2}}-i\sqrt{\frac{-a+\sqrt{a^2+b^2}}{2}} & b<0.
\end{array}
\right. \nonumber 
\end{eqnarray}
\end{lemma}

\begin{lemma}\label{lemma47}
\begin{itemize}
\item[(i)]When $\lambda \in \Xi_{+}$, $\frac{\omega}{\lambda} \left( \frac{-\lambda + \lambda^{-1} + \sqrt{\lambda^2 + \lambda^{-2}}}{\sqrt{2}} \right) = \frac{-1 \pm i}{\sqrt{2}}$ if and only if $\lambda = R_{-}(\omega) \mp i\mathrm{sgn}(\omega) R_{+}(\omega)$.
\item[(ii)]When $\lambda \in \Xi _{-}$, $\frac{\lambda}{\omega}(\frac{-\lambda +\lambda^{-1}+\sqrt{\lambda^2 +\lambda^{-2}}}{\sqrt{2}})=\frac{1\pm i}{\sqrt{2}}$ if and only if $\lambda = -R_{-}(\omega) \pm i \mathrm{sgn}(\omega) R_{+}(\omega)$.
\end{itemize}
\end{lemma}

\begin{proof}
($\mathrm{i}$) We solve the following equation with respect to $\lambda$:
\begin{eqnarray*}
\frac{\omega}{\lambda} \left( \frac{-\lambda + \lambda^{-1} + \sqrt{\lambda^2 + \lambda^{-2}}}{\sqrt{2}} \right) = \frac{-1 \pm i}{\sqrt{2}}. 
\end{eqnarray*}
By computing and simplifying this equation, we get
\begin{align}
    \label{413}
    \sqrt{\lambda^2+\lambda^{-2}}=(-1\pm i)\cdot \frac{\lambda}{\omega}+\lambda-\lambda^{-1}.
\end{align}
Squaring both sides of this equation and simplifying, we obtain
\begin{align*}
\lambda^2 &= \frac{(\omega^2 - \omega \pm i\omega)(-\omega \mp i(\omega -1))}{2\omega^2 -2\omega +1} \\
&= \frac{\omega (-\omega ^2 +2\omega -1) \mp i\omega (\omega ^2 -\omega +1)}{2\omega ^2 -2\omega +1}. 
 \end{align*}
We divide the problem into two cases. We first consider the case of 
\begin{align*}
\lambda ^2= \frac{\omega (-\omega ^2 +2\omega -1) - i\omega (\omega ^2 -\omega +1)}{2\omega ^2 -2\omega +1}. 
 \end{align*}
Since $\omega^2-\omega+1$ is positive for any $\omega \in \mathbb{R} \backslash \{0\}$, we observe that
\begin{eqnarray*}
\Im \lambda^2 = \left\{
\begin{array}{ll}
{\rm positive} &  {\rm if} \quad  \omega <0,\\
{\rm negative} & {\rm if} \quad \omega >0.
\end{array}
\right. 
\end{eqnarray*}
By Lemma \ref{lemma46} and the fact that $\Re \lambda \geq 0$,
\begin{eqnarray*}
\lambda = \left\{
\begin{array}{ll}
R_{-}(\omega)-iR_{+}(\omega) &   \omega >0,\\
R_{-}(\omega)+iR_{+}(\omega) &  \omega <0.
\end{array}
\right.
\end{eqnarray*}
The obtained expression for $\lambda$ can be written in the following form using $\mathrm{sgn}(\omega)$:
\begin{align*}
    \lambda = R_{-}(\omega)-i\mathrm{sgn}(\omega)R_{+}(\omega).
\end{align*}

 We next consider the case of 
\begin{align*}
\lambda ^2= \frac{\omega (-\omega ^2 +2\omega -1) + i\omega (\omega ^2 -\omega +1)}{2\omega ^2 -2\omega +1}, 
 \end{align*}
where 
\begin{eqnarray*}
\Im \lambda^2 = \left\{
\begin{array}{ll}
{\rm positive} &  {\rm if} \quad  \omega >0,\\
{\rm negative} & {\rm if} \quad \omega <0.
\end{array}
\right.
\end{eqnarray*}
Similarly as above, we obtain
\begin{align*}
    \lambda = R_{-}(\omega)+i\mathrm{sgn}(\omega)R_{+}(\omega).
\end{align*}

We need to verify whether the obtained values of $\lambda$ realy satisfy Eq.(\ref{413}), 
because we square Eq.\eqref{413} to solve the equation. Since $\Re \sqrt{\lambda^2 +\lambda^{-2}} \ge 0$, it is enough to see that $\Re \Bigl( (-1\pm i) \cdot \frac{\lambda}{\omega}+\lambda -\lambda^{-1} \Bigr) > 0$ for $\lambda = R_-(\omega) \mp i\mathrm{sgn}(\omega) R_+(\omega)$. 
Substituting $\lambda$ yields 
\begin{align*}
    (-1 \pm i) \cdot \frac{\lambda}{\omega} +\lambda -\lambda^{-1} &= -\frac{R_-(\omega)\mp i\mathrm{sgn}(\omega)R_+(\omega)}{\omega} + \frac{\pm iR_-(\omega) + \mathrm{sgn}(\omega)R_+(\omega)}{\omega}\\
    &+R_-(\omega)\mp i\mathrm{sgn}(\omega)R_+(\omega)-\frac{R_-(\omega)\pm i\mathrm{sgn}(\omega)R_+(\omega)}{R_-(\omega)^2 +R_+(\omega)^2}.
\end{align*}
By extracting the real part from both sides, we have
\begin{align*}
    \Re \Bigl( (-1\pm i) \cdot \frac{\lambda}{\omega} +\lambda -\lambda^{-1} \Bigr) =-\frac{R_-(\omega)}{\omega}+\frac{\mathrm{sgn}(\omega)R_+(\omega)}{\omega}+R_-(\omega) -\frac{R_-(\omega)}{R_-(\omega)^2+R_+(\omega)^2}.
\end{align*}
Considering $R_-(\omega) >0$, we only need to show
\begin{align}
    -\frac{1}{\omega}+ \mathrm{sgn}(\omega)\frac{1}{\omega}\frac{R_+(\omega)}{R_-(\omega)}+1 - \frac{1}{R_-(\omega)^2+R_+(\omega)^2} > 0. \label{5236}
\end{align}
Since $(\omega -1)^4+(\omega^2-\omega+1)^2=(2\omega^2-2\omega+1)(\omega^2-2\omega+2)$,
\begin{align*}
    \frac{R_+(\omega)}{R_-(\omega)} = \frac{\mathrm{sgn}(\omega)(\omega -1)^2 +\sqrt{(2\omega ^2 -2\omega +1)(\omega ^2 -2\omega +2)}}{\omega^2 -\omega +1}.
\end{align*}
Note that $2\omega^2 - 2\omega + 1 > 0$ and $\omega^2 - 2\omega + 2 > 0$ for any $\omega \in \mathbb{R} \setminus \{0\}$. Then, we find
\begin{align*}
R_-(\omega)^2+R_+(\omega)^2 &= \mathrm{sgn}(\omega) \omega \sqrt{\frac{\omega^2-2\omega +2}{2\omega^2 -2\omega +1}},
\end{align*}
and hence,
\begin{align*}
    \frac{1}{R_-(\omega)^2+R_+(\omega)^2} =\mathrm{sgn}(\omega)\frac{1}{\omega}\sqrt{\frac{2\omega^2 -2\omega +1}{\omega^2 -2\omega +2}}.
\end{align*}
Substituting these results into Eq.(\ref{5236}) and
%\begin{align*}
%-\frac{1}{\omega} + \mathrm{sgn}(\omega) \frac{1}{\omega}\frac{\mathrm{sgn}(\omega)(\omega -1)^2+\sqrt{(2\omega^2-2\omega+1)(\omega^2-2\omega+1)}}{\omega^2-\omega+1} + 1-\mathrm{sgn}(\omega)\frac{1}{\omega}\sqrt{\frac{2\omega^2-2\omega+1}{\omega^2-2\omega+2}} > 0.
%\end{align*}
multiplying by $\mathrm{sgn}(\omega) \omega(\omega^2 - \omega + 1)\sqrt{\omega^2-2\omega +2}$, we obtain that Eq.\eqref{5236} is equivalent to
\begin{align}
\label{4716}
(\omega-1)(\mathrm{sgn}(\omega)\omega^2 \sqrt{\omega^2-2\omega +2} -\sqrt{2\omega^2-2\omega+1})>0.
\end{align}

When $\omega > 0$, the equation
\begin{align*}
    \left(\omega^2 \sqrt{\omega^2 - 2\omega + 2}\right)^2 - \left(\sqrt{2\omega^2 - 2\omega + 1}\right)^2 =(\omega-1)(\omega^2+1)(\omega^2-\omega+1),
\end{align*}
shows that $\omega^2 \sqrt{\omega^2-2\omega +2} -\sqrt{2\omega^2-2\omega+1}$
is positive when $\omega>1$ and is negative when $0<\omega<1$.
Therefore, Eq.(\ref{4716}) is true.
When $\omega < 0$, it is clear that Eq.(\ref{4716}) holds.


($\mathrm{ii}$) We solve the equation
\begin{eqnarray*}
\frac{\lambda}{\omega}\Bigl(\frac{-\lambda +\lambda^{-1}+\sqrt{\lambda^2 +\lambda^{-2}}}{\sqrt{2}}\Bigr)=\frac{1\pm i}{\sqrt{2}}
\end{eqnarray*}
with respect to $\lambda$.
Rewriting the equation yields the following expression:
\begin{align}
    \label{416}
    \sqrt{\lambda^2+\lambda^{-2}}=(1\pm i)\cdot \frac{\omega}{\lambda}+\lambda-\lambda^{-1}.
\end{align}
Squaring both sides of this equation and simplifying, we obtain
\begin{align}
\lambda^2 &= \frac{(\pm i\omega^2 - \omega \mp i\omega)(1-\omega \pm i\omega )}{2\omega^2 -2\omega +1} \nonumber \\
&= \frac{\omega (-\omega ^2 +2\omega -1) \mp i\omega (\omega ^2 -\omega +1)}{2\omega ^2 -2\omega +1}. \nonumber
 \end{align}
We divide the problem into two case. We first consider the case of 
\begin{align*}
\lambda ^2= \frac{\omega (-\omega ^2 +2\omega -1) - i\omega (\omega ^2 -\omega +1)}{2\omega ^2 -2\omega +1},
 \end{align*}
where
\begin{eqnarray*}
\Im \lambda^2 = \left\{
\begin{array}{ll}
{\rm positive} &  {\rm if} \quad  \omega <0,\\
{\rm negative} & {\rm if} \quad \omega >0.
\end{array}
\right. 
\end{eqnarray*}
By Lemma \ref{lemma46} and the fact that $\Re \lambda \leq 0$,
\begin{eqnarray*}
\lambda = \left\{
\begin{array}{ll}
-R_{-}(\omega)+iR_{+}(\omega) &   \omega >0,\\
-R_{-}(\omega)-iR_{+}(\omega) &  \omega <0.
\end{array}
\right.
\end{eqnarray*}
It follows from the above that $\lambda$ can be written as
\begin{align*}
    \lambda = -R_{-}(\omega)+i\mathrm{sgn}(\omega)R_{+}(\omega).
\end{align*}

We next consider the case of 
\begin{align*}
\lambda ^2= \frac{\omega (-\omega ^2 +2\omega -1) + i\omega (\omega ^2 -\omega +1)}{2\omega ^2 -2\omega +1}, 
 \end{align*}
where 
\begin{eqnarray*}
\Im \lambda^2 = \left\{
\begin{array}{ll}
{\rm positive} &  {\rm if} \quad  \omega >0,\\
{\rm negative} & {\rm if} \quad \omega <0.
\end{array}
\right.
\end{eqnarray*}
Similarly as above, we obtain
\begin{align*}
    \lambda = -R_{-}(\omega)-i\mathrm{sgn}(\omega)R_{+}(\omega).
\end{align*}
We now verify that the obtained values of $\lambda$ realy satisfy Eq.(\ref{416}). Since $\Re \sqrt{\lambda^2 +\lambda^{-2}} \ge 0$, it is enough to show that $\Re \Bigl( (1\pm i) \cdot \frac{\lambda}{\omega}+\lambda -\lambda^{-1} \Bigr) > 0$ for $\lambda = -R_-(\omega) \pm i\mathrm{sgn}(\omega) R_+(\omega)$.
Substituting $\lambda$ yields
\begin{align*}
    (1 \pm i) \cdot \frac{\omega}{\lambda} +\lambda -\lambda^{-1} =& -R_{-}(\omega)\pm i\mathrm{sgn}(\omega)R_{+}(\omega)+(-1+\omega)\cdot \frac{-R_{-}(\omega)\mp i\mathrm{sgn}(\omega)R_{+}(\omega)}{R_-(\omega)^2+R_+(\omega)^2}\\
    &+i\omega\frac{\mp R_{-}(\omega)-i\mathrm{sgn}(\omega)R_{+}(\omega)}{R_-(\omega)^2+R_+(\omega)^2}.
\end{align*}
By extracting the real part from both sides, we get
\begin{align*}
    &\Re \Bigl( (1\pm i) \cdot \frac{\lambda}{\omega} +\lambda -\lambda^{-1} \Bigr) =\\
    &-R_-(\omega)-(-1+\omega)\cdot \frac{R_-(\omega)}{R_-(\omega)^2+R_+(\omega)^2}+\mathrm{sgn}(\omega)\omega\cdot  \frac{R_+(\omega)}{R_-(\omega)^2+R_+(\omega)^2}.
\end{align*}
Since $R_-(\omega)>0$, we only need to prove
\begin{align}
    -1-(-1+\omega)\cdot \frac{1}{R_-(\omega)^2+R_+(\omega)^2}+\mathrm{sgn}(\omega)\omega \cdot \frac{R_+(\omega)}{R_-(\omega)}\cdot  \frac{1}{R_-(\omega)^2+R_+(\omega)^2} > 0. \label{4723}
\end{align}
Substituting the expressions for $\frac{R_+(\omega)}{R_-(\omega)}$ and $\frac{1}{R_-(\omega)^2 + R_+(\omega)^2}$ computed in (i) into Eq.(\ref{4723}) and multiplying by 
$\mathrm{sgn}(\omega)\omega(\omega^2 - \omega + 1)\sqrt{\omega^2-2\omega+2}$,
we obtain that Eq.\eqref{4723} is equivalent to 
\begin{align*}
(\omega-1)(\mathrm{sgn}(\omega)\omega^2 \sqrt{\omega^2-2\omega +2} -\sqrt{2\omega^2-2\omega+1})>0.
\end{align*}
The above inequality coincides with Eq.(\ref{4716}). Therefore, we have $\Re \Bigl( (1 \pm i) \cdot \frac{\lambda}{\omega} +\lambda -\lambda^{-1} \Bigr) > 0$.
\end{proof}

\end{document}
